\documentclass[10pt,a4paper]{amsart}
\usepackage[margin=19mm]{geometry}
\usepackage{amsmath,amssymb,mathtools}
\usepackage[scr=rsfs,cal=euler]{mathalfa}
\usepackage{xfrac}
\usepackage[unicode,hidelinks,bookmarks=false]{hyperref}
\newcommand{\ie}{i.e.\ }
\newcommand{\cf}{cf.\ }
\newcommand{\resp}{respectively\ }
\newcommand{\Subsection}{\subsection}
\providecommand{\nobreakdash}{\nobreak\hbox{-}}
\setlength{\emergencystretch}{2em}
\multlinegap0pt
\usepackage{amssymb}
\usepackage{mathtools}
\newtheorem{lemma}{Lemma}
\newtheorem{proposition}{Proposition}
\newcommand{\IKl}{\mathrm{IKl}}
\newcommand{\Nrd}{\mathrm{Nrd}}
\newcommand{\Trace}{\mathrm{Tr}}
\newcommand{\Trd}{\mathrm{Trd}}
\renewcommand{\geq}{\geqslant}
\renewcommand{\subset}{\subseteq}
\title{Exotic and inverted Kloosterman sums over semisimple algebras}
\author{Daqing Wan, Dingxin Zhang}
\date{Source: arXiv:2606.05771v1 (CC BY 4.0)}
\begin{document}
\maketitle

\noindent\emph{Source and licence.} Exotic and inverted Kloosterman sums over semisimple algebras. Version arXiv:2606.05771v1 (CC BY 4.0), distributed by arXiv under the Creative Commons Attribution 4.0 International licence, http://creativecommons.org/licenses/by/4.0/, as recorded in the arXiv licence metadata for this exact version and shown on the abstract page https://arxiv.org/abs/2606.05771v1. Full licence notice: \texttt{licenses/arxiv-2606.05771-CC-BY-4.0.txt}.\par\smallskip

\noindent\textit{Editorial scope.} Counted unit: the article's proposition prop:inverted-master (source0.tex lines 632--643). The article's complete original proof of it is delivered with it (source0.tex lines 645--788, one \texttt{proof} environment of 144 lines). No mathematics is written, repaired or re-derived: the statement and the proof are the article's own lines, in the article's own notation, and no retained line is substituted or re-flowed.  Retained uncounted as the source-local context the proof needs: lines 407--419; lines 606--617.  Nothing was omitted from any retained span, so the package carries no omission record.  No retained line contains a citation, so the wrapper carries no bibliography and no bibliography entry.  No retained line contains a figure environment, an image-inclusion command or drawing code, so no image is delivered and the package declares no asset dependency.  Editorial formatting only, in addition: an AMS \texttt{amsart} preamble that restates the article's own declarations and macros used by the retained lines, namely the environments \texttt{lemma}, \texttt{proposition} on separate counters, together with the standard packages those lines need.  The retained environments print in this document as Lemma 1, Proposition 1 in order of appearance, whereas the article prints the same items with the numbering of its own full document.  The complete native source of the article as distributed in the arXiv e-print for this exact version is bundled unchanged as \texttt{sources/202140/source0.tex}.
\begin{lemma}\label{lem:matrix-reduction}
Let \(F=\mathbb{F}_{q^d}\), let
\(\psi_F=\psi\circ\Trace_{F/\mathbb{F}_q}\), and let
\(Q=q^d\).  For \(e\geq 1\) and \(y\in F^\times\),
\begin{equation}\label{eq:matrix-reduction}
\sum_{\substack{A\in\mathrm{GL}_e(F)\\ \det A=y}}
\psi_F(\operatorname{Tr}A)
=
Q^{\binom{e}{2}}
\sum_{\substack{x_1,\ldots,x_e\in F^\times\\  x_1\cdots x_e=y}}
\psi_F(x_1+\cdots+x_e).
\end{equation}
\end{lemma}

\begin{lemma}\label{lem:inverted-affine}
Let \(L\colon\mathbb{A}^r(\mathbb{F}_q)\to\mathbb{F}_q\) be a
nonconstant affine-linear function.  Then
\[
\sum_{\substack{z\in\mathbb{A}^r(\mathbb{F}_q)\\L(z)\neq 0}}
\psi(L(z)^{-1})
=
-q^{r-1}.
\]
This is the inverted-character analogue of the cancellation identity
\(\sum_{z\in\mathbb{F}_q}\psi(cz)=0\) for \(c\neq 0\).
\end{lemma}

\begin{proposition}\label{prop:inverted-master}
With notation as above, for every \(a\in\mathbb{F}_q^\times\),
\begin{equation}\label{eq:inverted-master}
\IKl_M(a;\chi)
=
q^N\, \IKl_{B^{\prime}}(a;\eta)
-
\left(\sum_{\substack{w\in W\\ w\neq 1}} q^{N-1+\ell_d(w)}\right)
\sum_{\substack{t\in (B^{\prime})^\times\\ \nu(t)=a}} \eta(t).
\end{equation}
The sum over \(w\neq 1\) is interpreted as zero if \(W=\{1\}\).
\end{proposition}

\begin{proof}
Write \(F_i=\mathbb{F}_{q^{d_i}}\).  In the \(i\)-th factor
\(\mathrm{GL}_{n_i}(F_i)\), let \(\mathcal{B}_i\) be the group of
invertible upper triangular matrices, and let \(\mathcal{U}_i\) be the
subgroup of upper triangular matrices with diagonal entries all equal
to \(1\).

For each \(w_i\in S_{n_i}\), choose a determinant-one monomial
representative \(\dot w_i\), as in the proof of
Lemma~\ref{lem:matrix-reduction}.  For
\(w=(w_1,\ldots,w_k)\in W\), put
\(\dot w=(\dot w_1,\ldots,\dot w_k)\).  Also put
\(\mathcal{B}_M=\prod_i\mathcal{B}_i(F_i)\).  The product of the
matrix cell decompositions is the disjoint union
\[
M^\times
=
\bigsqcup_{w\in W}
\mathcal{U}_w\, \dot w\, \mathcal{B}_M,
\]
where
\[
\mathcal{U}_w=
\mathcal{U}_{w_1}(F_1)\times\cdots\times
\mathcal{U}_{w_k}(F_k).
\]
Here \(\mathcal{U}_{w_i}(F_i)\subset \mathcal{U}_i(F_i)\) is the
small affine space from Lemma~\ref{lem:matrix-reduction}: for
\(w_i=(w_i(1),\ldots,w_i(n_i))\), one allows a free upper-triangular
entry \(u_{w_i(b),w_i(a)}\) for each inversion \(a<b\) with
\(w_i(a)>w_i(b)\), and sets all other upper off-diagonal entries to
\(0\).  Thus
\[
\#\mathcal{U}_w=q^{\sum_i d_i\ell(w_i)}=q^{\ell_d(w)}.
\]

Decompose the inverted sum accordingly:
\begin{equation}\label{eq:ikl-cell-decomp}
\IKl_M(a;\chi)=\sum_{w\in W}\IKl_w(a;\chi),
\end{equation}
where
\[
\IKl_w(a;\chi)
=
\sum_{\substack{u^{\prime}\in \mathcal{U}_w,\, b\in \mathcal{B}_M\\
    \Nrd(u^{\prime}\dot w b)=a\\ \Trd(u^{\prime}\dot w b)\neq 0}}
\chi(b)\psi(\Trd(u^{\prime}\dot w b)^{-1}).
\]
The determinant-type character is unchanged by \(u^{\prime}\) and by the
determinant-one representative \(\dot w\).

Fix \(w\).  Since the representatives \(\dot w_i\) have determinant
\(1\), the condition \(\Nrd(u^{\prime}\dot w b)=a\) is simply
\(\Nrd(b)=a\).  Also, by cyclicity of the matrix trace in each factor,
\(\Trd(u^{\prime}\dot w b)=\Trd(bu^{\prime}\dot w)\).  For fixed
\(u^{\prime}\), the change of variables \(c=bu^{\prime}\) preserves
\(\mathcal{B}_M\), the reduced norm, and the character.  Hence
\[
\IKl_w(a;\chi)
=
\#\mathcal{U}_w
\sum_{\substack{c\in \mathcal{B}_M\\
    \Nrd(c)=a\\ \Trd(c\dot w)\neq 0}}
\chi(c)\psi(\Trd(c\dot w)^{-1}).
\]
Every upper triangular matrix \(c\in \mathcal{B}_M\) decomposes uniquely
as \(c=tu\), where \(t\in (B^{\prime})^\times\) is viewed as a diagonal
matrix and \(u\in \mathcal{U}\) is upper-unitriangular.  Here
\[
\mathcal{U}=\mathcal{U}_1(F_1)\times\cdots\times \mathcal{U}_k(F_k)
\]
is the full product of upper-unitriangular groups, an affine space of
dimension \(N=\sum_i d_i\binom{n_i}{2}\) over \(\mathbb{F}_q\).  Using
cyclicity once more, \(\Trd(tu\dot w)=\Trd(u\dot w t)\).  Substituting
\(c=tu\) gives
\begin{equation}\label{eq:ikw-expanded}
\IKl_w(a;\chi)
=
q^{\ell_d(w)}
\sum_{\substack{t\in (B^{\prime})^\times\\ \nu(t)=a}}
\eta(t)
\sum_{\substack{u\in \mathcal{U}\\
    \Trd(u\dot w t)\neq 0}}
\psi(\Trd(u\dot w t)^{-1}).
\end{equation}

When \(w=1\), \(\Trd(ut)=\Trd(t)\) for every \(u\in \mathcal{U}\).
Thus \eqref{eq:ikw-expanded} gives
\begin{equation}
\label{eq:w1}
\IKl_1(a;\chi)
=
q^N\IKl_{B^{\prime}}(a;\eta).
\end{equation}

If instead \(w\neq 1\), then some \(w_i\) is a non-identity permutation,
so its monomial representative has a nonzero entry below the diagonal.
Say this entry is in row \(s\), column \(r\), with \(s>r\).  In the
factor \(\mathcal{U}_i\), the coordinate \(u_{rs}\) is a free variable.
In this factor, put \(Y_i=\dot w_i t_i\).  Then
\[
\operatorname{Tr}(u_iY_i)
=
\operatorname{Tr}(Y_i)+\sum_{a<b}(u_i)_{ab}(Y_i)_{ba},
\]
so \(u_{rs}\) occurs in \(\operatorname{Tr}(u_i\dot w_i t_i)\) with
coefficient \((Y_i)_{sr}\neq 0\).  Here the corresponding entry of
\(\dot w_i\) is nonzero, and every diagonal entry of \(t_i\) is nonzero.

After applying the field trace
\(\Trace_{F_i/\mathbb{F}_q}\), the resulting function of the
\(F_i\)-variable \(u_{rs}\) is still nonconstant, because the pairing
\((c,z)\mapsto \Trace_{F_i/\mathbb{F}_q}(cz)\) is nondegenerate.  Thus
\[
u\longmapsto \Trd(u\dot w t)
\]
is a nonconstant affine-linear function on the \(N\)-dimensional
\(\mathbb{F}_q\)-affine space \(\mathcal{U}\).  Lemma~\ref{lem:inverted-affine}
therefore gives
\[
\sum_{\substack{u\in \mathcal{U}\\ \Trd(u\dot w t)\neq 0}}
\psi(\Trd(u\dot w t)^{-1})
=
-q^{N-1}.
\]
Substituting in \eqref{eq:ikw-expanded} yields, for \(w\neq 1\),
\begin{equation}
\label{eq:w2}
\IKl_w(a;\chi)
=
-q^{\ell_d(w)}q^{N-1}
\sum_{\substack{t\in (B^{\prime})^\times\\ \nu(t)=a}}\eta(t).
\end{equation}

Adding these contributions \eqref{eq:w1},\eqref{eq:w2} in
\eqref{eq:ikl-cell-decomp} gives
\[
\IKl_M(a;\chi)
= q^N \, \IKl_{B^{\prime}}(a;\eta)
- \sum_{\substack{w\in W\\w\neq 1}} q^{N-1+\ell_d(w)}
\sum_{\substack{t\in (B^{\prime})^\times \\ \nu(t)=a}}\eta(t).
\]
This is \eqref{eq:inverted-master}.
\end{proof}

\end{document}
